\begin{center}\begin{tikzpicture}[font=\small,>=stealth,line join=round]
\draw[gray!40,rounded corners=3pt] (0,0) rectangle (14.2,8.3);\node at (7.1,7.700000000000001) {例6：改变系数，伸缩或翻转负指数的曲线};
\begin{scope}[shift={(1.6853658536585368,3.9499999999999997)},x=1.4634146341463412cm,y=0.2421052631578947cm]\draw[->,gray] (-0.4,0)--(3.7,0) node[right]{$x$};\draw[->,gray] (0,-9.5)--(0,9.5) node[above]{$y$};\draw[blue,thick,solid,domain=0.6:3.2,samples=180] plot(\x,{1/((\x)^2)});\draw[teal,thick,dashed,domain=0.6:3.2,samples=180] plot(\x,{3/((\x)^2)});\draw[orange!85!black,thick,densely dotted,domain=0.6:3.2,samples=180] plot(\x,{-3/((\x)^2)});\fill[orange!85!black] (1,1) circle(2.2pt);\fill[orange!85!black] (1,3) circle(2.2pt);\fill[orange!85!black] (1,-3) circle(2.2pt);\node[above right=5pt] at (1,3){$(1,3)$};\end{scope}\node[anchor=west,align=left,text width=5.0cm] at (8.4,4.0) {蓝：$x^{-2}$\\[8pt]青：$3x^{-2}$，仍递减\\[8pt]橙：$-3x^{-2}$，变递增\\[8pt]$(1,3)$确定$A=3$\\[8pt]后两者非严格定义的幂函数};
\node[text width=12.6cm,align=center] at (7.1,.65) {同一个坐标系，只有竖直高度被系数改变；负系数还把图像翻到横轴下方。};
\end{tikzpicture}\end{center}
